%===============================================================================
\section{Introduction}\label{sec:intro}
%===============================================================================

Electricity markets publish a large and growing amount of operational data. Disclosure rules
list, field by field, what must be released and when: forecasts and requirements before
clearing, prices and their components right after clearing, and actual operating quantities on
the following day~\cite{nea2024disclosure}. Other quantities---the output or the offers of an
individual generating unit, the hourly flow on a specific line, or the actual load before it is
disclosed---are protected, either permanently or until a later release date. Data
classification standards ask operators to grade each field and, explicitly, to consider how
aggregation and processing change a field's level~\cite{gbt43697}. How to do the latter is
left open.

The difficulty is that the inference risk of a field is not a property of the field alone. Power
data come from one physical and economic system, and the information a field adds depends on what
else is known. A unit's hourly output says little about its private offer, and neither does the
price at its bus; together, whenever the unit is marginal, they reveal the offer
(Section~\ref{sec:res-mech}). The same price adds nothing over one background and a great deal over
another. This is the inference and aggregation problem of database
security~\cite{lunt1989aggregation,farkas2002inference}, and it is not hypothetical for power data:
public prices alone can expose grid topology~\cite{kekatos2014topology}. Recent work on data
sensitivity assessment has recognized that a field's risk depends on its inference relations to
sensitive data and propagates risk along an inference graph~\cite{hu2026ignn}. What is still missing
is a field-level risk that is \emph{defined} by what an adversary can infer from a set of
fields---so that it can be checked by an attack experiment rather than calibrated to expert
scores---and that captures inference requiring several fields jointly.

We start from the gain $\Delta_i(T)$ that releasing field $i$ brings over a background $T$ of fields
the adversary already holds, measured as the out-of-sample accuracy of retrained attackers. A field
then has a whole distribution of contextual gains, and grading it means summarizing that
distribution. Attribution methods such as Shapley values and SAGE average the gains over backgrounds;
this is the right answer to ``how much does the field contribute typically'', but a security audit
asks whether \emph{some} admissible background makes the field dangerous. We therefore take the
maximum---the Marginal Contribution Index~\cite{catav2021mci}, in line with worst-case measures of
leakage in quantitative information flow~\cite{smith2009qif,alvim2012gleakage}---and add the two
conditions a disclosure audit needs: the data already public and a bound on the adversary's
background. The resulting budgeted marginal risk $\Mik$, together with its single-field part
$M^{(0)}$, the \emph{background amplification} $\Gamma^{(K)}=\Mik-M^{(0)}$ and the top backgrounds
that attain it, is the field-level output of the audit.

Computing it requires the inferability of every small field set, and it has to be recomputed
whenever targets or public data change. We amortize this with one network pre-trained under random
visibility masks---an adversary that does not know in advance which fields it will obtain---and a
closed-form readout solved per field set; the estimate is only used to scan, and a handful of
backgrounds per field is retrained to certify the result. We found that pre-training on the
sensitive targets alone contributes almost nothing beyond random features, because the representation
collapses onto a few target directions; our estimator therefore combines masked reconstruction of all columns with the
features of an untrained network and a safeguard against exploding predictions.

\paragraph{Contributions}
\begin{itemize}
\item \emph{Background-aware field risk.} From set-level inferability we define the budgeted
marginal risk $\Mik$ and the background amplification $\Gamma^{(K)}$, relate them to the MCI, give a
closed form of the amplification at $K=1$, and show that $\Mik$ is the minimal threshold-free safety
margin and that its critical fields are exactly the members of small minimal unsafe sets
(Section~\ref{sec:mk}).
\item \emph{An amortized attacker with scan--certify.} Masked pre-training plus a per-set closed-form readout
estimates $\Vy$ for every small set in milliseconds; an error analysis explains why field risk is estimated more
accurately than the set values, certification by retraining gives lower bounds that do not depend on the estimate, and a
diagnosis of the pre-training shows why a target-only objective collapses and motivates the design of the estimator
(Sections~\ref{sec:model} and~\ref{sec:res-est}).
\item \emph{Mechanism and real-data evidence.} A controlled power-system case with a known inference
mechanism, and four data sets with rule-based field roles (RTS-GMLC, the Australian NEM with actual
unit outputs, PJM, CAISO) with $1.2\times10^5$ exhaustively retrained field sets as ground truth. We show
that background amplification is common and physically interpretable, compare the definition with
correlation-, attribution- and inference-graph-based scores, ablate the estimator, and test
robustness and the choice of $K$ (Sections~\ref{sec:res-mech}--\ref{sec:res-robust}); disposal is
discussed as an application (Section~\ref{sec:application}).
\end{itemize}

\begin{figure}[!tbp]
\centering
\includegraphics[width=\textwidth]{fig1_framework.pdf}
\caption{Audit framework. Disclosure rules fix field roles. Set-level inferability $V_y(S)$ is
defined by dedicated retraining and scanned by an amortized attacker; the scan proposes, for
every field, backgrounds that are certified by retraining. The field-level profile---single-field
risk, background amplification and top backgrounds---supports grading and downstream disposal.}
\label{fig:framework}
\end{figure}
